\section{Metodologia}\label{sec:Method}

Esta pesquisa é quantitativa, aplicada e retrospectiva, e utiliza dados secundários de partidas já realizadas \cite{wazlawick2014,marconiLakatos2021}. O procedimento foi organizado em cinco etapas: (i) coleta das fontes de dados; (ii) padronização e integração das fontes em uma base única de partidas; (iii) resolução de divergências e de valores ausentes; (iv) construção das variáveis preditoras a partir apenas de informações anteriores a cada partida; e (v) treinamento e avaliação dos modelos com validação cronológica.

\subsection{Escopo e unidade de análise}\label{subsec:Escopo}

O estudo abrange as partidas da Série A do Campeonato Brasileiro de 2017 a 2024, totalizando oito temporadas de 380 partidas, ou 3.040 partidas. A unidade de análise é a partida, identificada pela chave composta por data, equipe mandante e equipe visitante. A temporada de 2020, disputada de 8 de agosto de 2020 a 25 de fevereiro de 2021 em razão da pandemia de covid-19, foi tratada como uma única temporada: as 112 partidas realizadas entre 6 de janeiro e 25 de fevereiro de 2021 foram atribuídas a ela.

Os alvos são duas contagens por partida, modeladas separadamente: o total de gols (soma dos gols do mandante e do visitante) e o total de escanteios (soma dos escanteios das duas equipes). O resultado da partida (vitória do mandante, empate ou vitória do visitante) é derivado do placar e não é utilizado como variável preditora.

\subsection{Fontes de dados}\label{subsec:Fontes}

A base principal é o Brasileirão Dataset \cite{brasileiraoDatasetGithub}, que reúne, em arquivos separados, as informações de cada partida (data, hora, rodada, equipes, técnicos, formações e placar) e as estatísticas por equipe (chutes, chutes no alvo, posse de bola, faltas, cartões e escanteios). Essa fonte cobre as 3.040 partidas do período, mas 405 delas não possuem estatísticas registradas: toda a temporada de 2024, 23 partidas de 2019, uma de 2017 e uma de 2020. Por isso, outras cinco fontes foram utilizadas como validadoras e como origem de valores ausentes, conforme a Tabela~\ref{tab:fontes}. Essas fontes não são conferências independentes entre si: a base principal e o Sofascore, em especial, parecem ter a mesma origem (Seção~\ref{subsec:Limitacoes}).

\begin{table}[htb]
\IBGEtab{%
  \caption{Fontes de dados utilizadas}\label{tab:fontes}
}{%
  \begin{tabular}{p{3.4cm}p{2.4cm}p{6.6cm}}
    \toprule
    Fonte & Temporadas & Uso no estudo \\
    \midrule
    Brasileirão Dataset & 2017--2024 & Base principal: partidas, placares, técnicos, formações e estatísticas \\
    Sofascore & 2017--2024 & Segunda fonte na hierarquia: placares, estatísticas e horário de início; Copa do Brasil de 2017 a 2024 \\
    Kaggle 2014--2023 & 2017--2023 & Validadora de placares e escanteios \\
    Kaggle 2023--2024 & 2023--2024\footnotemark & Validadora de placares e estatísticas \\
    FBref & 2019, 2021--2024 & Validadora de placares e de escanteios por partida \\
    Football-Data & 2017--2024 & Validadora de placares \\
    \bottomrule
  \end{tabular}%
}{%
  \fonte{Elaborada pelo autor (2026).}
}
\end{table}
\footnotetext{Partidas de 2024 até 28 de agosto.}

As fontes são o Brasileirão Dataset \cite{brasileiraoDatasetGithub}, o Sofascore \cite{sofascore}, o Brazilian Football Matches \cite{brazilianFootballMatches}, aqui chamado Kaggle 2014--2023, o conjunto Brasileirao \cite{kaggleBrasileirao2324}, aqui chamado Kaggle 2023--2024, o FBref \cite{fbref} e o Football-Data \cite{footballData}. A ordem de prioridade entre as fontes, usada para preencher ausências e desempatar divergências, é: Brasileirão Dataset, Sofascore, Kaggle e, por fim, FBref e Football-Data. Para os escanteios, a ordem é Brasileirão Dataset, Sofascore, FBref e Kaggle: o FBref passa à frente do Kaggle porque seus escanteios coincidiram com os da base principal em 98,9\% a 99,5\% das partidas de cada temporada e com os do Sofascore em 98,2\% a 100\%, contra 96,8\% a 98,9\% de coincidência entre o FBref e o Kaggle 2014--2023.

Além da Série A, foram reunidas as partidas da Copa do Brasil, usadas apenas para construir uma variável de calendário (Seção~\ref{subsec:Features}) e não incorporadas à base de partidas. Elas foram coletadas do Sofascore para todo o período, de 2017 a 2024. O Kaggle 2014--2023 também traz a Copa do Brasil, mas não foi usado para esse fim, porque lhe faltam partidas, entre elas jogos de ida ou de volta de confrontos entre equipes da Série A.

\subsection{Coleta}\label{subsec:Coleta}

As fontes disponíveis como arquivos (Brasileirão Dataset, os dois conjuntos do Kaggle, o calendário do FBref e o Football-Data) foram obtidas diretamente de seus repositórios. Duas fontes exigiram coleta própria. Do FBref, foram extraídos os relatórios de cada partida, dos quais se aproveitou a tabela de estatísticas complementares, totalizando 1.900 partidas das temporadas de 2019 e de 2021 a 2024, com escanteios em 1.885 delas; para 2017, 2018 e 2020, a fonte não disponibiliza relatórios de partida. As faltas do FBref não foram usadas, pois coincidiram com as da base principal em apenas 93,8\% das partidas comparáveis.

Os dados do Sofascore foram coletados pelo autor e armazenados em formato JSON sem alteração. Para cada partida, foram guardados o placar, as equipes, o horário de início e as estatísticas consolidadas do jogo e de cada tempo. A coleta reuniu as 3.040 partidas do período, 380 por temporada, todas com estatísticas. Da mesma forma, foram coletadas as partidas da Copa do Brasil de 2017 a 2024 que envolveram equipes que disputaram a Série A no período; nelas, o placar considerado é o do tempo regulamentar e da prorrogação, sem a disputa por pênaltis.

\subsection{Padronização e integração}\label{subsec:Integracao}

Cada fonte foi carregada no banco de dados PostgreSQL sem tratamento, com todas as colunas como texto, preservando os valores originais para auditoria. Em seguida, foram aplicados filtros para manter apenas a Série A no período do estudo, e cada fonte foi convertida para um esquema comum de colunas.

A padronização envolveu três tratamentos. O primeiro foi o dos nomes das equipes: as grafias de cada fonte foram mapeadas manualmente para um nome único por clube (177 pares de fonte e grafia para 31 clubes nas seis fontes da Série A), inclusive no caso de clubes que mudaram de denominação no período. O segundo foi o das datas: os horários do Kaggle 2014--2023 e do Sofascore estão em tempo universal coordenado (UTC), e os do Football-Data, no horário de Londres; todos foram convertidos para o horário de Brasília antes da extração da data, para que partidas noturnas não fossem atribuídas ao dia seguinte. O terceiro foi o das unidades: posse de bola foi registrada como inteiro de 0 a 100, e, no Kaggle 2023--2024, o total de chutes foi obtido pela soma dos chutes fora do alvo e no alvo, e a posse registrada como 0 a 0 foi tratada como ausente.

As fontes foram então unidas à base principal pela chave da partida. A base principal define quais partidas existem; partidas presentes só nas validadoras não entram no estudo. Após a padronização, 2.649 das 2.651 partidas do Kaggle 2014--2023 no período e todas as partidas do Kaggle 2023--2024 (610), do FBref (1.900) e do Football-Data (3.040) encontraram correspondência na base principal. As duas exceções diferem apenas na data registrada.

Alguns campos de contexto da base principal foram corrigidos antes da integração. O horário de início foi substituído pelo do Sofascore nas 49 partidas em que as duas fontes diferiam em mais de 15 minutos, o que inclui 8 das 13 partidas registradas antes das 11 horas na base principal. Formações e técnicos ausentes ou atípicos em 23 partidas foram conferidos em súmulas da Confederação Brasileira de Futebol (CBF) e em portais esportivos (ESPN, GE e G1). Em 25 partidas do Paraná em 2018, a base principal registra como técnico o auxiliar Ademir Fesan, que dirigiu a equipe em apenas uma delas; nessas partidas, o técnico foi corrigido para Rogério Micale, até sua demissão em 14 de agosto, e para Claudinei Oliveira, a partir daí \cite{tribunaMicale2018,tribunaClaudinei2018}.

\subsection{Divergências e valores ausentes}\label{subsec:Conflitos}

A ausência de um valor não foi tratada como zero. Na base principal, as 405 partidas sem estatísticas aparecem com estatísticas iguais a zero e posse nula; esses valores foram tratados como ausentes. Quando uma estatística não estava disponível em uma fonte, ela foi mantida como ausente e, se possível, recuperada de outra fonte, com o registro da origem de cada valor de gols e de escanteios na base final.

As divergências entre fontes foram tratadas por um protocolo aplicado na seguinte ordem:
\begin{enumerate}
  \item zero suspeito: um zero na base principal é tratado como ausente quando alguma validadora registra valor positivo para a mesma partida;
  \item ausência: um valor ausente na base principal é preenchido pela primeira fonte da hierarquia que o possua;
  \item divergência sistemática: para cada atributo e temporada, uma validadora que diverge da base principal em mais de 10\% das partidas comparáveis deixa de ser usada naquele atributo e temporada;
  \item gols divergentes: prevalece o valor da maioria das fontes disponíveis, com desempate pela hierarquia;
  \item demais atributos divergentes: regras específicas por atributo, descritas a seguir.
\end{enumerate}

Considera-se divergência qualquer diferença em gols e escanteios, diferença maior que um chute no total de chutes e diferença maior que um ponto percentual na posse de bola. A comparação entre as fontes mostrou concordância de 100\% nos gols entre a base principal, o Sofascore, o FBref e o Football-Data. Nos escanteios, a base principal coincidiu com o Sofascore em 100\% das partidas comparáveis, com o FBref em 99,1\% e com o Kaggle em 99,0\%. Nos chutes, o Kaggle 2014--2023 divergiu da base principal acima do limite de 10\% e foi excluído desse atributo. O Kaggle 2023--2024 também divergiu acima do limite (17,7\% dos chutes) e, por exceção à regra, foi mantido, porque só preenche um valor da base final (os chutes da partida conferida manualmente, descrita adiante).

As regras por atributo foram as seguintes. Nos gols, houve seis divergências, todas resolvidas pela maioria das fontes, que confirmou o valor da base principal. Nos escanteios, quando a base principal e o Sofascore concordam, prevalece esse valor, mesmo que o FBref ou o Kaggle divirjam, o que ocorreu em 38 partidas. Sem valor na base principal, prevalece a maioria entre Sofascore, FBref e Kaggle e, sem maioria, o Sofascore; se este também estiver ausente, vale a hierarquia. Nos chutes, chutes no alvo, posse e faltas divergentes, prevalece a base principal e, na sua falta, o Sofascore. Ao todo, 449 casos foram registrados: seis de gols e 443 de escanteios, dos quais 393 foram preenchimentos de ausência, 38 resolvidos pela base principal, seis pela maioria, cinco pela hierarquia e um por conferência manual. Este último é a partida Athletico-PR x Internacional, de 21 de abril de 2024, em que as duas únicas fontes com escanteios divergiam (1 a 1 no FBref e 5 a 3 no Kaggle 2023--2024) e a hierarquia indicava o FBref; a conferência manual, feita pelo autor nas estatísticas exibidas pelo Google na busca da partida, consultadas em 3 de novembro de 2026, confirmou 5 a 3.

A base final tem 3.040 partidas, 380 por temporada, com chave única e sem nenhum valor ausente nos gols e nos escanteios. Para cada partida, a base registra a fonte que originou os gols e os escanteios, e o resultado da partida é recalculado a partir do placar final.

\subsection{Variáveis preditoras e prevenção de vazamento}\label{subsec:Features}

A previsão é feita antes do início da partida, depois da divulgação das escalações. Por isso, os gols, escanteios, chutes, posse de bola, faltas e cartões da partida-alvo não são usados como entrada do modelo: essas estatísticas só existem depois do jogo. Elas entram apenas como histórico, calculado sobre partidas anteriores de cada equipe. As partidas de cada equipe são ordenadas por data e hora, e não pela rodada, porque 75 partidas foram adiadas e disputadas fora da ordem das rodadas; o histórico é deslocado em uma partida antes de qualquer agregação, de modo que a estatística da própria partida nunca entra no seu cálculo. Médias da liga, valores de referência e eventuais transformações das variáveis são calculados apenas com dados anteriores ou ajustados somente com as partidas de treino.

O primeiro grupo de variáveis é a média, na temporada corrente, de gols e escanteios a favor e contra de cada equipe. No início da temporada, essa média se apoia em poucas partidas e é instável; por isso, ela é encolhida em direção a um valor de referência, conforme a Equação~(\ref{eq:enc}):
\begin{equation}\label{eq:enc}
  \tilde{x} = \frac{n\,\bar{x} + k\,r}{n + k},
\end{equation}
em que $\bar{x}$ é a média da equipe nas $n$ partidas já disputadas na temporada, $r$ é o valor de referência e $k$ controla o peso da referência, que diminui à medida que $n$ cresce. A referência é a média final da equipe na temporada anterior. Para as equipes promovidas da Série B, cuja temporada anterior não está na base, a referência é a média das equipes promovidas nas temporadas anteriores do estudo e, na falta destas, a média da liga na temporada anterior. Um indicador de equipe promovida também entra como variável. A temporada de 2017, primeira do período, não tem temporada anterior; por isso, as dez partidas da sua primeira rodada ficam sem histórico, e o indicador de promovida não é definido nesse ano. O valor de $k$ é escolhido na validação cronológica entre 0, 2, 5, 10, 20, 40, 80, 160, 320 e 1.000 (Seção~\ref{subsec:Modelos}); com $k = 0$, não há encolhimento, e com $k = 1.000$ a média fica, na prática, igual à referência.

O segundo grupo é formado por médias móveis das últimas $w$ partidas, com $w = 5$ e $w = 10$, em dois recortes: geral, com as partidas em casa e fora, e por mando, em que o mandante é descrito só pelas partidas em casa e o visitante só pelas partidas fora. As médias abrangem gols, escanteios, chutes e chutes no alvo, a favor e contra, e a posse de bola a favor, o que resulta em 72 variáveis. As janelas não atravessam temporadas; nas primeiras rodadas, com menos de $w$ partidas disponíveis, aplica-se o mesmo encolhimento da Equação~(\ref{eq:enc}).

Como alternativa às janelas de tamanho fixo, as mesmas estatísticas, nos mesmos recortes, são resumidas por médias com decaimento exponencial: todas as partidas anteriores da temporada contam, mas o peso de cada uma cai à metade a cada $h$ partidas. São calculadas meias-vidas de 3, 5 e 10 partidas, e a média é sempre encolhida pela Equação~(\ref{eq:enc}), com $n$ igual ao número de partidas já disputadas na temporada. A escolha entre janelas e decaimento, e o valor de $h$, são feitos na validação cronológica.

O terceiro grupo resume a força das equipes por \textit{pi-ratings}, na versão simplificada de \citeonline{constantinou2022}, derivada de \citeonline{constantinouFenton2013}. Cada equipe tem um índice para as partidas em casa e outro para as partidas fora, na unidade de saldo de gols. Na partida entre a equipe $X$, mandante, e a equipe $Y$, visitante, o saldo esperado é a diferença entre o índice de $X$ em casa e o índice de $Y$ fora. Depois do jogo, o índice de $X$ em casa aumenta em $\lambda e$, em que $e$ é o saldo observado menos o esperado, e o índice de $X$ fora recebe a fração $\gamma$ dessa variação; os índices de $Y$ variam no sentido oposto, da mesma forma. Usam-se as taxas do artigo original, $\lambda = 0{,}035$ e $\gamma = 0{,}7$. Os índices de uma partida são lidos antes da atualização com o seu placar. A temporada de 2017 começa com todos os índices em zero. Nas seguintes, cada equipe começa com uma fração $f$ do índice final da temporada anterior, e o restante é levado a zero, a média da liga, para refletir as mudanças de elenco entre temporadas; as equipes promovidas começam com a mesma fração da média dos índices finais das promovidas em temporadas anteriores. O valor de $f$ é escolhido na validação cronológica entre 0,5, 0,7 e 1.

Como os \textit{pi-ratings} descrevem apenas a diferença entre as equipes, o total de gols e de escanteios é resumido também pelos índices GAP de \citeonline{wheatcroft2020}. Cada equipe tem quatro índices, de ataque e de defesa, em casa e fora, para cada medida de desempenho ofensivo: gols, escanteios, chutes, chutes no alvo e a soma de chutes e escanteios. A produção esperada do mandante é a média entre o seu índice de ataque em casa e o índice de defesa do visitante fora, e a do visitante, a média entre o seu índice de ataque fora e o índice de defesa do mandante em casa; a soma das duas estima o total da partida. Depois do jogo, os índices são revistos pela diferença entre o observado e o esperado, conforme as equações do artigo original, e não podem ficar negativos. As três taxas de cada medida foram escolhidas por busca em grade apenas com a temporada de 2017, que não entra na avaliação. Uma equipe que aparece pela primeira vez começa na média da liga das partidas com início anterior, e na virada da temporada aplica-se a mesma regra dos \textit{pi-ratings}, com a média final da liga no lugar de zero.

O quarto grupo descreve o calendário e o contexto da partida, ambos conhecidos antes do jogo. Os dias de descanso medem o intervalo entre o início da partida anterior da equipe, na Série A ou na Copa do Brasil, e o início da partida-alvo, considerando apenas partidas da Copa efetivamente realizadas. O valor é limitado a 14 dias, próximo do percentil 99 dos intervalos dentro da temporada; intervalos maiores, como o que antecede a primeira partida da temporada, são tratados como descanso pleno, e a variável nunca fica ausente. Um indicador marca as partidas disputadas sem público durante a pandemia: toda a temporada de 2020 e as partidas de 2021 anteriores a 2 de outubro, quando a Série A voltou a receber torcedores, a partir da 23ª rodada, por decisão do Conselho Técnico da competição em 28 de setembro de 2021 \cite{opovoPublico2021}.

O técnico da partida-alvo também é conhecido antes do jogo. A base principal grafa os nomes de forma abreviada e inconsistente, com variantes da mesma pessoa e pessoas distintas com a mesma grafia. Por isso, cada par de técnico e clube foi associado a uma pessoa pelo código do Sofascore, aceito apenas quando o nome confere com o da base principal; os nomes civis usados nessa conferência foram verificados em fontes públicas, e perfis duplicados do Sofascore foram unificados em um único identificador. Dos 395 pares, 352 receberam identificador; os 43 restantes, dos quais 38 têm até cinco partidas, são identificados pela grafia da base principal. Um técnico diferente em uma única partida, como em suspensões ou interinidades curtas, é tratado como o técnico anterior; a partir de duas partidas, o interino conta como técnico. Por isso, a troca só passa a contar na segunda partida seguida do novo técnico: antes dela, não se sabe se a substituição dura mais de uma partida. O técnico não entra como categoria, o que exigiria estimar um efeito para cada técnico, muitos com poucas partidas; entram o número de trocas de técnico da equipe na temporada e o número de partidas seguidas da equipe sob o técnico atual. A formação da partida-alvo é a registrada pela base principal como inicial; admite-se que ela seja conhecida antes do jogo, com a escalação divulgada, mas o momento em que a fonte define o desenho tático não foi verificado. As formações são agrupadas, por critério tático e não pelo número de escanteios, em quatro famílias conforme a amplitude pelos lados do campo: alta, com alas ou com laterais e meias abertos (por exemplo, 3-5-2, 4-4-2 e 4-5-1); média, com pontas e laterais mais fixos (4-3-3, 4-2-3-1 e 4-1-4-1); estreita, sem pontas nem alas (por exemplo, 4-1-2-1-2 e 4-2-2-2); e com cinco defensores (5-4-1 e 5-3-2). Uma segunda variável conta os jogadores do meio-campo para a frente. Uma terceira resume os escanteios a favor e contra da equipe nas partidas anteriores em que usou a mesma família, em todas as temporadas, no geral e só no mesmo mando, encolhidos pela Equação~(\ref{eq:enc}) para a média da liga na família, calculada apenas com partidas de início anterior. Os modelos são avaliados com e sem a formação, e o ganho é medido, sobretudo para os escanteios.

Ficam fora do estudo os cartões amarelos e vermelhos, o público presente e o resultado da partida; do público, entra apenas o indicador de partida sem torcida. O conjunto de variáveis usado por modelo é descrito na Seção~\ref{subsec:Modelos}.

\subsection{Modelos e avaliação}\label{subsec:Modelos}

Para cada alvo, o total de gols e o total de escanteios da partida, são comparadas quatro previsões da média. A referência ingênua é a média do alvo nas temporadas de treino, a mesma para todas as partidas. O modelo linear é a regressão de Poisson com penalização \textit{ridge} \cite{hoerlKennard1970}, ajustada sobre as variáveis padronizadas pela média e pelo desvio-padrão do treino; a penalização reduz a variância das estimativas quando há variáveis correlacionadas, como as várias medidas de força das equipes. Os dois modelos não lineares são o Random Forest \cite{breiman2001}, com 300 árvores e critério de divisão pela desviância de Poisson, e o XGBoost \cite{chenGuestrin2016}, com função objetivo de Poisson. A dispersão é medida no Poisson sem penalização, pela estatística de Pearson dividida pelos graus de liberdade; acima de cerca de 1,1, ajusta-se também a regressão binomial negativa (NB2), com o parâmetro de dispersão $\alpha$ estimado pela regressão auxiliar de \citeonline{cameronTrivedi1990}.

As variáveis entram em um conjunto enxuto por alvo. As médias a favor e contra viram o valor esperado do confronto: para o mandante, a média entre o que ele produz e o que o visitante sofre; para o visitante, o inverso. De cada grupo de variáveis entra uma única versão, escolhida na validação interna: a forma recente (janelas de 5 ou 10 partidas ou decaimento com meia-vida de 3, 5 ou 10, no recorte geral ou por mando), a fração $f$ dos \textit{pi-ratings}, a fração dos índices GAP e, só para os escanteios, a formação (ausente, pelas famílias de amplitude ou pelo número de jogadores do meio-campo para a frente), com os escanteios por família apenas no recorte por mando. Os índices GAP entram apenas com as medidas do próprio alvo. Completam o conjunto os dias de descanso, o indicador de partida sem público e as trocas de técnico na temporada; a posse de bola fica de fora. O indicador de equipe promovida e o número de partidas sob o técnico têm ausências documentadas e entram apenas no XGBoost, que trata valores ausentes na construção das árvores. O conjunto tem 16 variáveis nos gols e até 24 nos escanteios; no XGBoost, com as duas variáveis com ausências de cada equipe, são 20 e até 28.

A avaliação respeita a ordem temporal das partidas, em validação progressiva por temporada (\textit{walk-forward}). Para cada temporada de teste $T$, de 2019 a 2024, os modelos são treinados com as temporadas anteriores a $T$ e testados em $T$. Nunca se usa divisão aleatória, que permitiria ao modelo aprender com jogos futuros. As temporadas de 2017 e 2018 entram apenas no treino: 2017 como aquecimento das variáveis e 2018 como primeira temporada de validação interna.

Todas as escolhas de uma dobra usam apenas as temporadas anteriores a $T$. Cada temporada de validação $v$, de 2018 a $T-1$, é prevista por um modelo ajustado nas temporadas anteriores a $v$, e cada configuração candidata recebe a desviância média dessas validações. A configuração com a menor média é reajustada com todas as temporadas anteriores a $T$ e prevê $T$. A média de várias validações foi adotada porque a escolha por uma única temporada mostrou-se instável entre dobras vizinhas. No modelo linear, a grade combina $k$ (Equação~\ref{eq:enc}), a forma recente, as frações dos \textit{pi-ratings} e dos índices GAP, a formação e a penalização, com 14 valores entre $10^{-6}$ e 3.000. As árvores usam o $k$ e as variáveis escolhidos pelo modelo linear na mesma dobra, e a busca se concentra nos hiperparâmetros. No Random Forest, a grade completa combina o número mínimo de partidas por folha (de 10 a 300), a fração de variáveis sorteadas em cada divisão (de 0,1 a 1,0) e a profundidade máxima (1, 2, 3, 4, 6 ou sem limite), em 252 combinações. No XGBoost, sorteiam-se, sem repetição, 240 das 15.552 combinações, 120 com cada uma de duas sementes fixas, de profundidade máxima, peso mínimo por folha, frações de partidas e de variáveis por árvore, penalizações L1 e L2 e taxa de aprendizado. O número de árvores, de 5 a 1.500, é escolhido pelas previsões parciais do mesmo ajuste, sem parada antecipada pela temporada de teste. As grades de $k$ e do Random Forest foram ampliadas em duas etapas, sempre que a escolha caiu na borda. Nas bordas restaram apenas limites naturais, como $k = 1.000$, com o qual a média da temporada se reduz à referência, a profundidade 1 e a fração 1,0 de variáveis. No XGBoost, a amostra foi ampliada uma vez, e a borda da taxa de aprendizado (0,1) foi mantida, porque o número de árvores, escolhido na mesma validação, compensa a taxa. Uma ampliação adicional da busca do XGBoost, feita como análise de sensibilidade, é discutida na Seção~\ref{sec:Discussion}.

Avaliam-se também duas combinações das previsões. A primeira é a média ponderada das médias previstas pela referência ingênua, pelo modelo linear, pelo Random Forest e pelo XGBoost, com pesos não negativos que somam 1, buscados em uma grade de passo 0,1 e escolhidos em cada dobra pela mesma desviância média das temporadas de validação de 2018 a $T-1$, sem a temporada de teste. Não há meta-modelo ajustado sobre as previsões, apenas a escolha dos pesos. Como os hiperparâmetros de cada modelo foram escolhidos nessas mesmas validações, as previsões usadas para escolher os pesos favorecem um pouco os modelos ajustados. A segunda é a média simples do Random Forest e do XGBoost, sem escolha de pesos; ela foi definida depois de vistos os resultados de teste dos dois modelos nas rodadas anteriores e antes do resultado da própria média.

O desempenho foi medido pelo erro absoluto médio (MAE), pela raiz do erro quadrático médio (RMSE) e pela desviância de Poisson média, definidos nas Equações~(\ref{eq:mae}) a (\ref{eq:dev}), em que $y_i$ é a contagem observada na partida $i$, $\hat{y}_i$ é a contagem prevista e $n$ é o número de partidas avaliadas:
\begin{equation}\label{eq:mae}
  \mathrm{MAE} = \frac{1}{n}\sum_{i=1}^{n} \left| y_i - \hat{y}_i \right|,
\end{equation}
\begin{equation}\label{eq:rmse}
  \mathrm{RMSE} = \sqrt{\frac{1}{n}\sum_{i=1}^{n} (y_i - \hat{y}_i)^2},
\end{equation}
\begin{equation}\label{eq:dev}
  D = \frac{2}{n}\sum_{i=1}^{n} \left[ y_i \ln\!\left(\frac{y_i}{\hat{y}_i}\right) - (y_i - \hat{y}_i) \right],
\end{equation}
com $y_i \ln(y_i/\hat{y}_i) = 0$ quando $y_i = 0$. O MAE e o RMSE medem o erro na mesma unidade do alvo, sendo o RMSE mais sensível a erros grandes; a desviância de Poisson avalia a previsão como média de uma distribuição de contagens, o que é relevante para os gols, evento de baixa frequência, e para os escanteios, de frequência maior.

As previsões de cada modelo são comparadas com a referência ingênua pela diferença de desviância por partida, com intervalo de confiança de 95\% por \textit{bootstrap} com 2.000 reamostragens. As partidas de uma mesma rodada compartilham o calendário e a fase do campeonato e não são independentes; por isso, a reamostragem sorteia rodadas inteiras, identificadas pela temporada e pelo número da rodada, e não partidas isoladas, em um \textit{bootstrap} por blocos \cite{kunsch1989}. Uma diferença é tratada como significativa quando o intervalo não contém o zero. O intervalo mede a incerteza da média sobre as 2.280 partidas de teste, com os modelos já escolhidos e ajustados; ele não inclui a variação entre temporadas como unidade, a variação da escolha de hiperparâmetros nem a correção para comparações múltiplas.

Como análise de robustez, posterior à obtenção dos resultados e não planejada de início, os modelos também são comparados com duas referências mais fortes que a média da liga, no mesmo \textit{bootstrap}. A primeira é a média do total na temporada anterior, sem parâmetro algum. A segunda é a média corrente da liga: a Equação~(\ref{eq:enc}) com $k = 100$, em que $\bar{x}$ é a média do total nas $n$ partidas da temporada com data estritamente anterior à da partida (partidas do mesmo dia não entram) e $r$ é a média da temporada anterior. O peso de 100 partidas foi fixado de uma vez, sem ajuste, e depois de vistos os resultados de teste.

Além da média, avalia-se a probabilidade de o total passar de uma linha usual de apostas: mais de 2,5 gols e mais de 9,5 escanteios. Com a média prevista por cada modelo, a probabilidade é obtida pela distribuição de Poisson e pela binomial negativa. O parâmetro $\alpha$ desta última é estimado, em cada dobra, pela regressão auxiliar de \citeonline{cameronTrivedi1990} sobre as previsões fora da amostra das temporadas de validação interna, sem uso da temporada de teste. A referência é a frequência de partidas acima da linha nas temporadas de treino. As probabilidades são avaliadas pela perda logarítmica e pelo escore de \citeonline{brier1950}, com o mesmo \textit{bootstrap} por rodadas.

A contribuição das variáveis é medida por importância por permutação \cite{breiman2001}, na temporada de teste de cada dobra: cada variável, ou cada grupo de variáveis relacionadas, como as de formação, é embaralhada entre as partidas, sem reajustar o modelo, e registra-se o aumento médio da desviância em 20 repetições. A permutação por grupo evita que variáveis correlacionadas dividam a importância entre si. No XGBoost, calcula-se também a contribuição média absoluta de cada variável pelos valores SHAP \cite{lundbergLee2017}.

\subsection{Limitações dos dados}\label{subsec:Limitacoes}

Quatro limitações decorrem das fontes. Primeiro, a base principal e o Sofascore concordam em 100\% dos escanteios comparáveis, o que sugere origem comum dos dados; as duas fontes, portanto, não são validações independentes, e a confirmação de um valor por ambas não equivale a duas conferências. Segundo, os campeonatos estaduais, a Copa Libertadores e a Copa Sul-Americana não estão na base, de modo que a carga de partidas é subestimada e os dias de descanso, superestimados. Terceiro, os chutes vêm de mais de uma fonte, com definições que podem diferir, como a soma de finalizações no Kaggle 2023--2024; em 2024, as estatísticas vêm do Sofascore em 379 partidas e do Kaggle 2023--2024 em uma. Nos mesmos jogos de 2017 a 2023, porém, o Sofascore e a base principal coincidem em 99,3\% dos totais de chutes, o que reduz o risco de que a troca de fonte em 2024 desloque a distribuição. Quarto, os dados do Sofascore foram coletados pelo autor e o código dessa coleta não é distribuído; a base integrada final será disponibilizada com o código do trabalho \Rl{https://github.com/PedroVieiraC/previsao-gols-escanteios-brasileirao}, mas essa etapa não pode ser repetida por terceiros.
